% LSPV-4 REPORT::CLASSCHARGE1::2026-10-05T00:51:00-07:00
\documentclass[11pt]{article}
\usepackage[margin=.86in]{geometry}
\usepackage{amsmath,amssymb}
\begin{document}
\section*{LSPV-4 CLASSCHARGE1: hostile classification-cost audit}
\textbf{GLOBAL STATUS: OPEN.} P versus NP remains open. This is conditional implementation accounting only; it proves no SAT shrinker, acquisition theorem, causal compiler, circuit lower bound, or barrier-crossing hinge.
\section*{CORPUS REVIEWED}
Governance, CODEX, pnp-research skill, root main.tex, MAIN2303, BATCHSETUP1, SHORTKERN1, SOURCEEDGE1, LOCALBIT1, TASK-008/slice lineage, and HEC/CRS-PPC/CCSPC search surfaces were reviewed. LOCALBIT1 verifies as predecessor. Direct fetch and repository search find no BATCHFIX1 artifact; MAIN2303 and SOURCEEDGE1 independently record that discrepancy, so its fallback formula is not promoted.
\section*{Primary-literature preflight}
Razborov--Rudich applies only when constructivity, largeness, and usefulness hypotheses are met under its pseudorandomness assumption. Williams sharpens the constructive-property/derandomization relation; almost-natural work shows weakening largeness matters. The 2026 switching-lemma barrier is for an $AC^0$-natural subclass, not a generic TASK-008 refutation.
\section*{PROVED/EXACT conditional classification accounting}
Let $Q=n^{\eta+o(1)}$ required children fall into $D=n^{\delta+o(1)}$ source-proved reusable SAME\_REALIZATION setup classes. Setup costs $n^{p+o(1)}$, generation $n^{g+o(1)}$, child scale is $n^{\theta+o(1)}$, and child exponent is $C_j$.
If each child computes its class at cost $n^{\kappa_{\rm ch}+o(1)}$, then
\[\boxed{C_i\ge\max\{r_i,\eta+\kappa_{\rm ch},\delta+p,\eta+g+\theta C_j\}.}\]
If each class is enumerated at cost $n^{\kappa_{\rm cl}+o(1)}$ and setup then costs $n^p$, sequential overhead is
\[n^\delta(n^{\kappa_{\rm cl}}+n^p)=n^{\delta+\max\{\kappa_{\rm cl},p\}+o(1)},\]
hence
\[\boxed{C_i\ge\max\{r_i,\delta+\max\{\kappa_{\rm cl},p\},\eta+g+\theta C_j\}.}\]
The stronger $\delta+p+\kappa_{\rm cl}$ requires an actual multiplicative $n^{p+\kappa_{\rm cl}}$ per-class computation. Sequential polynomial costs do not generically add exponents.
Implementation-independently, if total classification costs $K_{\rm class}=n^{\chi+o(1)}$, use
\[\boxed{C_i\ge\max\{r_i,\chi,\delta+p,\eta+g+\theta C_j\}.}\]
If classification, setup, generator, and child workspace exponents are $b,a,s,\rho_j$, class-by-class depth-first execution gives
\[\boxed{\rho_i\ge\max\{\sigma_i,b,a,s,\theta\rho_j\}.}\]
\section*{Hostile controls}
RECOMP trades state for repeated classification time. INTERP compresses representation, not evaluation. DELAY requires realized rather than semantic class counts. COPY can raise $\eta$ without raising $\delta$ or $\chi$. A current-scale call to hypothetical SAT decider $A$ inside classification reintroduces its unknown exponent $K$. Semantic/algebraic oracles may collapse classification/setup, so this theorem relativizes.
\section*{Structural and causal audit}
SHORTKERN1 and LOCALBIT1 remain structural/source-realization results only. LOCALBIT1 explicitly retires automatic promotion of $p^r\to p^{r/2}$ to SAT $\theta=1/2$. No source-certified numerical tuple $(\eta,\theta,\delta,p,\chi,a,b,g,s)$ exists. Classification also does not repair SOURCESTOP1/PATHSTRAT1/OFFPATH1: compact off-path strategies do not force arbitrary static SAT deciders to acquire counterfactual information.
\section*{TASK-008 and Natural-Proofs audit}
CLASSCHARGE1 gives no TASK-008 lower bound. DISTEQUIV1/Valiant keep the exact-slice target calibrated as general-circuit-strength after transfer. Fixed-order camouflage remains a hostile filter. Broad promoted invariants must face polynomial-size unrestricted distinguishers and, where applicable, Razborov--Rudich hypotheses. The 2026 switching-lemma barrier remains model-specific.
\section*{DUPLICATION / CONTRADICTION AUDIT}
BATCHSETUP1 says classification must be separately charged but leaves its formula unspecified. The fallback BATCHFIX1 term $\delta+p+\kappa$ is not generically implied. Per-class sequential classification gives $\delta+\max\{\kappa,p\}$; per-child classification gives $\eta+\kappa_{\rm ch}$. No BATCHFIX1 artifact is currently readable.
\section*{WHAT CHANGED RELATIVE TO ROOT MAIN}
Root retains $\delta_e+p_e+\kappa_e$ as conservative candidate accounting while noting BATCHFIX1 is unrecovered. This term needs a semantic definition of $\kappa_e$. The implementation-independent replacement is separate total classification exponent $\chi_e$. No proof-spine arrow closes.
\section*{ADOPTED / REJECTED / RETIRED}
Adopt BATCHSETUP1, STREAMCOMP1, CYCLEGAIN1, SHORTKERN1, LOCALBIT1 at declared scope, SOURCESTOP1/PATHSTRAT1/OFFPATH1, DISTEQUIV1, SAME\_REALIZATION, and strict charge separation. Reject automatic exponent addition for sequential setup/classification and compact classifier description as free evaluation. Retire $\delta+p+\kappa$ as a generic ledger unless multiplicative semantics are explicit.
\section*{Single next target}
For the first actual source-matched Boundary-Squeeze edge, specify class identification as per-child, per-class, or global and prove $(\eta,\theta,\delta,p,\chi,a,b,g,s)$ with SAME\_REALIZATION grouping. Highest-value falsification is a hidden current-scale $A$ call or repeated RECOMP classification cost destroying the exponent.
\section*{References}
A. A. Razborov and S. Rudich, \emph{Natural Proofs}, JCSS 55(1):24--35, 1997. R. R. Williams, \emph{Natural Proofs versus Derandomization}, SIAM J. Comput. 45(2):497--529, 2016. T. Y. Chow, \emph{Almost-Natural Proofs}, FOCS 2008. B. Loff et al., \emph{The Switching Lemma shows what the Switching Lemma cannot prove}, 2026.
\bigskip\noindent\textbf{VERIFICATION SENTINEL: CLASSCHARGE1 complete; P versus NP OPEN.}
% END_OF_VERIFIED_REPORT::LSPV-4::CLASSCHARGE1
\end{document}